%% kapa: exam=1 round=33 year=2022 period=1 form=A subject=civil_law subject_ko=민법 seq=40 number=40
%% source: Q-Net 감정평가사 기출문제 (cst00309) artlSeq=5209997 file=raw/1008/5209997/2022년도 제33회 감정평가사 1차시험 1교시(민법,경제학,부동산학).pdf
%% conversion: kordoc 4.24.51 -> tools/gen1.py (auto); review=proofread-sonnet; flags=none
\begin{question}{40}
저당권에 관한 설명으로 옳지 않은 것은? (다툼이 있으면 판례에 따름)
\begin{choices}
  \item 저당권은 그 담보한 채권과 분리하여 타인에게 양도하거나 다른 채권의 담보로 하지 못한다.
  \item 저당물의 소유권을 취득한 제3자는 그 저당물에 관한 저당권 실행의 경매절차에서 경매인이 될 수 있다.
  \item 특별한 사정이 없는 한 건물에 대한 저당권의 효력은 그 건물에 종된 권리인 건물의 소유를 목적으로 하는 지상권에도 미친다.
  \item 전세권을 목적으로 한 저당권이 설정된 후 전세권이 존속기간 만료로 소멸된 경우, 저당권자는 전세금반환채권에 대하여 물상대위권을 행사할 수 있다.
  \item 저당목적물의 변형물인 물건에 대하여 이미 제3자가 압류하여 그 물건이 특정된 경우에도 저당권자는 스스로 이를 압류하여야 물상대위권을 행사할 수 있다.
\end{choices}
\end{question}
\answer{5}
